% Tapis de Quillen, with Quillen's offprint and preprint (cote 162-5, typed part): archive cote 162-5, site page 4 (PDF page 5).
% typescript. Draft compared: none.
% AI-transcribed by Claude Opus 5.5 (claude-opus-5-5, adaptive thinking, effort high), Message Batches API, batch t162_full (msgbatch_01UAh3zyWPd454AeDkhnmQ33), 2026-11-13; not yet checked by eye.
\scanpage{162-5/4}
\mspage{\hand{3}}

% NOTE: the right edge of the page is cut off on the scan; the last letters of most lines are not visible
premier. Comme application, Quillen obtient une interprétation faisceau-%
que de la cohomologie d'un ensemble semi-simplicial à coéfficients dans
système \inserted{local} \unsure{contravariant} \inserted{F} \xout{\unsure{de}}(défini classiquement par le complexe cosimplici-
\margin{$\Pi$ !} des $C^n(F) = \coprod_{x \in X_n} F(x)$ ): on considère le système local contravariant
défini par F, on l'interprète comme un faisceau sur \inserted[typed]{T(X) i.e. objet de} $(\mathrm{Ssimpl})_{/X}$ , et on
prend sa cohomologie. - Cependant, quand F est un système de coéfficients
\margin{[ (bracket in the left margin, from this line to ``faisceautique en termes classiques.'')} covariant pas nécessairement local, on n'a tjrs pas d'interprétation de
\margin{Si !} groupes de cohomologie classiques en t ermes \unsure{faiscdautiques|faisceautiques}; ni, lorsque
est \unsure{contravariant}, de son homologie, ou inversement de sa cohomologie
faisceautique en termes classiques.

A propos de la notion de foncteur qui est un homotopisme. Quillen
\margin{| (vertical line in the left margin, from this line to ``logue pour C'')} montre qu\xout{e}'un tel foncteur $F\colon C \to C'$ induit une équivalence entre la sous-
catégorie triangulée \inserted{$D^b_{lc}(C')$} de la catégorie \inserted{dérivée \inserted{bornée} de celle \deleted{des}} des faisceaux abéliens sur C', dont
les faisceaux de cohomologie sont des systèmes locaux, et la catégorie ana-
logue pour C; \xout{\unsure{Il existe}} et réciproquement. \xout{Pour la partie directe} On peut
dans cet énoncé introduire aussi n'importe quel anneau de base (à condition
de le supposer $\neq 0$ dans le cas de la réciproque); \xout{constant} la partie directe
vaut aussi avec un anneau de boefficients pas néc. constant, mais constant
tordu. Je pense que ce résultat (facile) doit pouvoir se généraliser ainsi
Soit $f\colon X \quad X'$ un morphisme de topos qui soit tel que pour tout \xout{systè}
faisceau loc. const. sur X', f induise un isomorphisme sur les cohomologies
(avec cas non commutatif inclus). Supposons que \deleted{\illegible{1} (X)} \inserted{\unsure{X} et X'} soit \uline{locale}
\margin{\o{} (call sign)} \margin{|| (double vertical line in the left margin, from this line to ``nuls pour'')} \uline{homotopiquement trivial}, i.e. que \xout{tout} pour tout entier n \inserted{$\geq 1$}, tout objet U de
\hand{\unsure{X'}} ait un \xout{\illegible{1}} \inserted[typed]{recouvrement} par des $U_i \to U$, tels que a) tout système local sur U
devient constant sur $U_i$ , et toute section sur U devient constante sur $U_i$
et b) pour tout \xout{système} groupe abélien G, les $H^j(U,G) \longrightarrow H^j(U_i,G)$ sont
nuls pour $1 \leq j \leq n$. Alors le foncteur $D^b_{lc}(X') \longrightarrow D^b_{lc}(X)$ induit par f est

\margin{\o{}} \hand{Attention, cette condition n'est \unsure{pratiquement} \unsure{pas} satisfaite par les schémas, avec leur topologie étale (mais bien par \deleted{\illegible{1}} \unsure{(noeth.} avec top. Zariski).}

\margin{4}
